\subsection{‘for’ Statements}\label{sec:ForStatements}
Since Java~5 the language knows two distinct types of \lstinline|for|\index{for} loops, both described in \autocite{ORACLE_DOC_LANGUAGE_SPECIFICATION:For}.

The basic \lstinline|for| statement\index{for!basic} looks like this:
\begin{lstlisting}[numbers=left]
for( <initialization>; <condition>; <update> )
{
    <statements>;
}

for( <initialization>; <condition>; <update> ) <singleStatement>;
\end{lstlisting}
For the details refer to \autocite{ORACLE_DOC_LANGUAGE_SPECIFICATION:ForBasic}.

An enhanced \lstinline|for| statement\index{for!enhanced} (defined in \autocite{ORACLE_DOC_LANGUAGE_SPECIFICATION:ForEnhanced}) would be written like below:

\begin{lstlisting}[numbers=left]
for( <declaration> : <iterable> )
{
    <statements>;
}

for( <declaration> : <iterable> ) <singleStatement>;
\end{lstlisting}

The use of blanks, round brackets, curly braces and linefeeds was already discussed in \tqfullvref{sec:WhiteSpace}.

Curly braces can be omitted only when there is just a single simple statement to be executed in the loop. In addition, this single statement has to be written completely into the same line as the \lstinline|for| itself. In any other case, the curly braces are mandatory.

An empty loop body can be written like this:
\begin{lstlisting}[numbers=left]
for( <initialization>; <condition>; <update> )
{
    // Deliberately empty
}

for( <initialization>; <condition>; <update> );

for( <declaration> : <iterable> )
{
    // Deliberately empty
}

for( <declaration> : <iterable> );
\end{lstlisting}

In case the longer forms are used, the comment is mandatory, because empty code blocks are forbidden. I recommend the shorter forms (lines~6 and 13).

When explicitly using an iterator\index{iterator}\index{java.util!Iterator}\autocite{ORACLE_DOC_ITERATOR_INTERFACE} with the (basic) \lstinline|for|\index{for!basic} loop, this should look like this:
\begin{lstlisting}
for( final var i = source.iterator(); i.hasNext(); )
{
    …
}
\end{lstlisting}

When using the comma operator in the initialization or update clause of a \lstinline|for|\index{for} statement, avoid the complexity of using more than three variables. If needed, use separate statements before the \lstinline|for|-loop (for the initialization clause) or at the end of the loop (for the update clause).

Examples:
\begin{lstlisting}
for( var i = 0, j = target.length; i < source.length && j >= 0; ++i, --j )
{
    target [j] = source [i];
}

var proceed = true;
var i = source.iterator();
for( var j = 0; proceed; )
{
    target [j++] = i.next();
    proceed = i.hasNext() && j < target.length;
}
\end{lstlisting}

\subsubsection{Principles}
\begin{enumerate}[label=P\arabic*.]
\item{one}
\item{two}
\item{three}
\end{enumerate}

\lstinline|for| loops are discussed in more detail in \tqfullvref{sec:TheForStatement}.

Examples:
\begin{lstlisting}[numbers=left]
// Ok
for( var i = 0; i < max; ++i ) sum += i;
for( final var s : texts ) System.out.println( s );

// DISCOURAGED
for( var i = 0; i < max; ++i ) if( v [i] > 0 ) sum += v [i];
for( final var s : texts ) if( !s.empty() ) out.println( s );

// AVOID!!!
for( var i = 0; i < max; ++i ) if( v [i] > 0 )
{
    sum += v [i];
}
else
{
    sum -= v [i];
}

for( final var s : texts ) if( !s.empty() )
{
    System.out.print( "Value: " );
    System.out.println( s );
}
\end{lstlisting}
The samples in lines~6 and 7 are not correct according to this rule, because the \lstinline|if| statement is not a simple statement..

An empty \lstinline|for|-loop (one in which all the work is done in the initialization, condition, and update clauses) should have the following form:\footnote{An \textit{enhanced} \lstinline|for|-loop without body does not make that much sense (that mentioned above is a classical \lstinline|for|-loop). It may be possible to write an implementation of \lstinline|java.lang.Iterable| with an iterator that does something as a side effect of \lstinline|hasNext()| or \lstinline|next()|, but this would break the contract of the interface \lstinline|java.lang.Iterator|.}
\begin{lstlisting}
for( <initialization>; <condition>; <update> );
\end{lstlisting}


%==============================================================================
% End of file!!
%==============================================================================
